\chapter{Brainteaser 完整答案}
\label{app:brainteasers-solutions}


\MFQInterviewFollowup{GB-2-01：若提案只需半数票且平票通过，逆推基例怎样改变？}
\MFQInterviewFollowup{GB-2-02：状态图出现平局循环时，胜负标记应增加哪一类状态？}
\MFQInterviewFollowup{GB-2-03：禁止空船返回后，如何用可达图证明无解？}
\MFQInterviewFollowup{GB-2-04：生日在二月二十九日前后会怎样改变星期增量？}
\MFQInterviewFollowup{GB-2-05：改变弃牌收益后，阈值策略为何可能移动或消失？}
\MFQInterviewFollowup{GB-2-06：若两根绳子的总燃烧时间不同，协议需新增什么校准信息？}
\MFQInterviewFollowup{GB-2-07：为什么信息量下界成立仍不保证称量设计可实现？}
\MFQInterviewFollowup{GB-2-08：在十二进制下，哪些素因子决定尾零数？}
\MFQInterviewFollowup{GB-2-09：若每轮只能比较四匹，候选淘汰表怎样重建？}
\MFQInterviewFollowup{GB-2-10：理论上消失的不稳定模态为何会被浮点误差重新激活？}
\MFQInterviewFollowup{GB-2-11：还需要哪些几何条件才能把必要条件加强为可行判定？}
\MFQInterviewFollowup{GB-2-12：若字体禁止把 $6$ 翻作 $9$，如何证明十二个面不足？}
\MFQInterviewFollowup{GB-2-13：若每个问题都必须预先写定而不能自适应，两问是否仍足够？}
\MFQInterviewFollowup{GB-2-14：有界重试能保证什么，又不能保证什么？}
\MFQInterviewFollowup{GB-2-15：改变同色与异色的放回规则后，应怎样重新寻找不变量？}
\MFQInterviewFollowup{GB-2-16：允许翻转任意距离的两个开关时，可达子空间如何变化？}
\MFQInterviewFollowup{GB-2-17：怎样用效用函数把奖金分布转成确定性等价？}
\MFQInterviewFollowup{GB-2-18：若只知道正面数的奇偶性而非精确值，原构造为何失效？}
\MFQInterviewFollowup{GB-2-19：若只有一个标签贴错，一次抽样是否仍足够？}
\MFQInterviewFollowup{GB-2-20：主持人的声明不是公共知识时，推理链断在哪一层？}
\MFQInterviewFollowup{GB-2-21：若允许曲线切割或切穿数字，整数下界还成立吗？}
\MFQInterviewFollowup{GB-2-22：缺失三个整数时还需要增加哪些幂和？}
\MFQInterviewFollowup{GB-2-23：测量误差如何限制可区分的二进制权重数？}
\MFQInterviewFollowup{GB-2-24：库存上界已知时，怎样求保证两种颜色各成一双的最坏抽数？}
\MFQInterviewFollowup{GB-2-25：有向握手或重复握手时，度数和公式怎样改写？}
\MFQInterviewFollowup{GB-2-26：把六人改成五人时，Ramsey 保证为何不再成立？}
\MFQInterviewFollowup{GB-2-27：蚂蚁速度不同后，标签交换等价的哪一步失效？}
\MFQInterviewFollowup{GB-2-28：称量有噪声时，编码之间需要保持怎样的最小距离？}
\MFQInterviewFollowup{GB-2-29：颜色数从二变为三时，一个公开比特为何不够？}
\MFQInterviewFollowup{GB-2-30：为什么换位错误通常无法由模 $9$ 校验发现？}
\MFQInterviewFollowup{GB-2-31：不变量给出不可达性后，怎样讨论其是否也是充分条件？}
\MFQInterviewFollowup{GB-2-32：若只知道正面数不超过 $n$，还能否保证两堆相等？}
\MFQInterviewFollowup{GB-2-33：允许叠放多块一次切开时，下界为什么改变？}
\MFQInterviewFollowup{GB-2-34：跑者速度随时间变化时，应积分哪个相对量？}
\MFQInterviewFollowup{GB-2-35：这是一种存在证明，为什么不必判定中间数是否有理？}
\MFQInterviewFollowup{GB-2-36：回答通道带噪声时，需要增加多少纠错冗余？}

\section*{逐题完整推导}
本节把每道题的“为什么”展开。先前的 36 条给出答案速览；这里补上状态、边界、构造和可复核的中间步骤。
\begin{enumerate}[leftmargin=*]
  \item 对剩余 $n$ 人定义状态 $V_n=(b_1,\ldots,b_n)$，其中每个分量是该人从当前轮开始的最优收益。基例 $V_1=(B)$；从 $V_{n-1}$ 反推时，当前提案者先把否决自己的结果写成 $V_{n-1}$，再按偏好从低到高购买最便宜的严格多数票，每票至少给下一轮所得加一，剩余奖金归自己。若平票通过，所需票数从 $\lfloor n/2\rfloor+1$ 改成 $\lceil n/2\rceil$，所以基例和每一层都必须重算；若总奖金不足，需声明平局偏好才能判断均衡。
  \item 有限状态写成 $(g,i_1,i_2,t)$，安全状态还要包含“谁行动”和是否已经捕获/逃脱。先把终止状态标为胜或负，再做吸引集迭代：对守卫回合，只要存在一条合法边进入守卫胜集就标胜；对闯入者回合，只有所有合法边都进入守卫胜集才标胜。循环未被标记时应另设 draw/unknown，不能把循环强行归为失败。
  \item 令左岸指示位为 $(r,m,k,d)$，安全约束是研究员不在某岸时，$K$ 不得与 $D$ 同岸，$M$ 不得与 $K$ 同岸。逐层 BFS 从 $(1,1,1,1)$ 出发，第一步只能带 $K$；回到左岸后带 $M$，再带 $K$ 回来，带 $D$ 过去，空返后带 $K$ 过去，共七条边到 $(0,0,0,0)$。删去任一空返边后，所有剩余节点都无法同时满足两项安全约束，图搜索直接证明无解。
  \item 把四年星期增量写成模 7 的差：平年同一日期前进 1，跨过二月二十九的区间前进 2。观测 $+1,+1,+2$ 只确定最后一个区间跨闰日；若生日与闰日的相对位置未给出，则“哪一个公历年份”仍有多个解。答案必须区分可确定的区间与不可确定的绝对日期。
  \item 设牌面为 $x$，加注的条件期望收益为 $U_R(x)$，弃牌/跟注基线为 $U_F(x)$。只有当 $U_R(x)-U_F(x)$ 随 $x$ 单调不减时，才能把策略写成阈值 $x\ge k$；阈值由底池、加注额、对手先验、弃牌收益和行动顺序共同决定。缺任一收益参数时只能证明候选形状，不能声称求出了均衡。
  \item 甲绳两端同时点燃，任意不均匀燃速仍在 30 分钟烧完；此时乙绳一端已烧 30 分钟，但剩余“燃烧时间”是 30 分钟而不是物理长度的一半。再点燃乙绳另一端，两个火头共同消耗这 30 分钟的燃烧量，故还需 15 分钟，总计 45 分钟。假设只需每端燃尽总时间稳定，不需均匀燃烧。
  \item 三次称量每次有左轻/平/右轻三类结果，信息量上最多 $3^3=27$ 条结果。异常状态是 12 个编号乘两种方向，共 24 个；因此三次是下界。构造时给每枚物品一个非零三元码 $c_i\in\{-1,0,1\}^3$，左/外/右对应分量，偏重输出 $c_i$、偏轻输出 $-c_i$；同时逐列检查左右盘数量相等、任何两枚码不相同也不相反。下界只说明可能，码表和每个结果的唯一解才完成构造。
  \item 用 Legendre 公式 $v_5(125!)=\lfloor125/5\rfloor+\lfloor125/25\rfloor+\lfloor125/125\rfloor=25+5+1=31$。十进制的 $2$ 因子更多，尾零数为 $\min(v_2,v_5)=31$；进制 $b=\prod p_j^{e_j}$ 时要计算 $\min_j\lfloor v_{p_j}(n!)/e_j\rfloor$，只数 5 只在十进制且 2 不成瓶颈时成立。
  \item 五组赛得到五个组冠军和组内名次；冠军赛确定 $A_1>B_1>C_1>D_1>E_1$。全局第一只能是 $A_1$，全局第二/三的候选为 $A_2,A_3,B_1,B_2,C_1$；第七轮让这五个候选赛一次，前两名就是全局第二、三。第六轮后至少仍有这五个候选的未比较排列，因此不可能保证前三，给出七轮下界。
  \item 特征方程 $r^2-3r+2=(r-1)(r-2)$，所以 $a_n=c_1+c_2 2^n$。由 $a_0=c_1+c_2$、$a_1=c_1+2c_2$ 解出 $c_2=a_1-a_0$；有限极限要求 $c_2=0$。浮点递推的误差也按 $2^n$ 放大，故应使用闭式或稳定重参数化并把误差增长作为输出。
  \item 单位立方体任意旋转后沿箱体最窄的 $0.9$ 方向投影宽度为 $|u_1|+|u_2|+|u_3|$，其中 $\|u\|_2=1$，故宽度至少为 1，大于 0.9，不能放入。体积和空间对角线只给必要条件；投影宽度是排除所有旋转的充分障碍。
  \item 日期 $01$--$09$ 要求两个方块都能显示 0、1、2 的配对；给出 $\{0,1,2,3,4,5\}$ 与 $\{0,1,2,6,7,8\}$ 后，6 翻转作为 9，即可覆盖 01--31。逐一检查个位/十位配对；若禁用 6/9 共用，面数约束给出不可能，说明字体约定是模型条件。
  \item 二元问题树深度 1 只有 2 个叶，不足以区分录用/加面/拒绝三类；信息下界为 $\lceil\log_2 3\rceil=2$。先问“是否录用”，否后再问“是否加面”，两个回答分支分别定位三类结果，达到下界。
  \item 若确认可能永久丢失，任意有限协议都有最后一条消息；它可能丢失而发送者不知道接收者是否知道，因此共同知识无法在有限轮形成。工程上可用序号、幂等、确认、超时和有界重试获得“去重+最终失败报告”，但不能同时保证必达和有限终止。
  \item 设当前蓝球数、总球数为 $(B_t,N_t)$。两蓝变一蓝使二者各减一；两红变一蓝使 $B$ 加一、$N$ 减一；异色变一红使二者各减一，故 $B_t+N_t\pmod2$ 不变。终局 $N_T=1$，若初始和奇则 $B_T=0$ 为红，若偶则 $B_T=1$ 为蓝。
  \item 状态向量在 $\mathbb F_2^n$ 中，翻转相邻两灯就是加 $e_i+e_{i+1}$。所有操作保持坐标和为 0；反过来从左到右，若当前位与目标不同就加 $e_i+e_{i+1}$，最后一位自动由偶校验决定。因此目标可达当且仅当亮灯奇偶相同，算法只需一次 $O(n)$ 扫描。
  \item 期望奖金只确定 $\mathbb E[B]$，效用比较还需分布尾部、支付时点、离职没收、税务流动性、与工资/公司收益相关性和风险厌恶。给两个相同均值但一个有大下行尾部的分布即可构成反例；没有效用函数和约束，只能比较期望现金流。
  \item 取任意 $k$ 枚为第一堆，若其中有 $h$ 枚正面，则第二堆有 $k-h$ 枚；翻转第一堆后正面数变为 $k-h$，与第二堆相等。关键是总正面数精确为 $k$；只知道奇偶性时无法写出 $k-h$，构造不成立。
  \item 从“混合”标签袋抽一次；标签全错意味着它不是混合，只能是纯 A 或纯 B，抽到的资产立即确定其真实类型。剩下两袋再用两个错误标签约束唯一交换，整个过程只需一次抽样。
  \item 主持人宣布至少一顶黑帽后删去全白世界。若第一人看见两白，他应立即知道自己黑，因此他说不知道排除相应世界；第二人把这一句话作为公共知识继续删世界；第三人的剩余集合只允许一种自身颜色。若声明不公开或不是公共知识，推理链不能继续。
  \item 四块若数字和相等，总和 $1+\cdots+12=78$ 必须被 4 整除，但 $78/4=19.5$ 非整数，因此不可能。这个结论依赖每个数字完整归属一块；允许切穿数字或按面积分配时模型改变。
  \item 缺失数 $x,y$ 满足 $s=x+y$、$q=x^2+y^2$，因此 $xy=(s^2-q)/2$。解二次式 $t^2-st+xy=0$，检查判别式为非负完全平方、两根为不同的 $1,\ldots,n$ 整数，并把补回后的总和/平方和重新核对；任一步失败都应拒绝输入。
  \item 为每组假币分配二进制权重 $2^i$，称量缺口除以单枚重量差后得到权重和，再由二进制展开恢复组集合。量程要容纳最大权重和，误差要小于相邻编码间隔一半；否则编码数学上可辨识但仪器上不可辨识。
  \item 保证一双同色需要 4 只（前三只各色，第四只必重复）。若要求两双且颜色不同，在每色库存无上限时不存在有限保证：可以先抽到任意多只同色，再各抽另一色一只；必须给出库存上界后才能做鸽巢计算。
  \item 把人为顶点、握手为无向边，每条边在两个端点度数中各计一次，故 $\sum_vd(v)=2|E|$。模 2 后奇数度顶点数为偶数；这不是经验规律，而是握手边的双端计数恒等式。
  \item 固定一人 $v$，其余五人中至少三人与 $v$ 的关系同类。若这三人中有一对认识，则与 $v$ 构成三人互识；若没有，则三人两两不认识。把认识与不认识交换，得到同样结论，这就是 $R(3,3)=6$ 的最小情形证明。
  \item 对每只蚂蚁只跟踪无标签轨迹。两只同速蚂蚁相撞掉头与彼此穿透、交换名字等价，因此无标签轨迹到出口的时间不变；第 $i$ 只沿初始方向到出口距离为 $d_i$ 时，总离开时间为 $\max_i d_i/v$。速度不同或碰撞有停顿时，交换标签不再等价。
  \item 三次称量码的三个分量分别为左/外/右。先枚举非零三元码的相反配对，再筛掉重复/相反码，最后检查每一列左右盘数量平衡；观察结果与码表匹配后唯一定位编号和轻重。若有噪声，还要给码之间留出 Hamming/加权距离，不能只用信息量下界。
  \item 第一人公开其余状态和的奇偶位；若总和为 $p$，他自己的位无法由自身观察恢复，可能成为牺牲者。之后第 $j$ 人用 $p$ 减去自己看到的其他位及此前已解出的位，在 $\mathbb F_2$ 中反解。保存的是一个全局校验位，不是每个人的独立答案。
  \item $N=\sum d_k10^k\equiv\sum d_k\pmod9$，因为 $10^k\equiv1$。单个数字错误若差不是 9 的倍数可检出；换位、多个误差抵消或 $0\leftrightarrow9$ 可能保持校验和，所以它只能检测一部分错误。
  \item 一次不同色相遇令两色各减一、第三色加二；模 3 等价于三个计数同时减一，两两差不变。初态 $8-4\equiv1$，全同色目标的差为 0，故不可达；不变量只给出必要障碍，若要讨论充分性还需证明所有同类状态之间可构造路径。
  \item 任取 $n$ 枚为第一堆，其中 $h$ 枚正面；总正面为 $n$，第二堆有 $n-h$ 枚，翻转第一堆后正面也为 $n-h$。这与 GB-2-18 同构，但这里 $n$ 同时是堆大小和总正面数；若只知道“不超过 $n$”，等式不再成立。
  \item 碎片数从 1 开始，每次掰一块变两块，净增 1；目标 $mn$ 块所以至少 $mn-1$ 次。沿网格线任意顺序掰开即可每次增加一块并达到下界。若允许叠放多块一次切开，单次增量不再是 1，原下界不适用。
  \item 同向相对速度为 $u-v$，第 $k$ 次追及满足 $(u-v)t_k=kL$，所以 $t_k=kL/(u-v)$，位置按 $ut_k\bmod L$。反向相对速度为 $u+v$，把分母替换为 $u+v$；必须先声明是否把 $t=0$ 的初始同点计为第 0 次。
  \item 令 $x=\sqrt2^{\sqrt2}$。若 $x$ 有理，取 $a=b=\sqrt2$，则 $a^b=x$ 有理；若 $x$ 无理，取 $a=x,b=\sqrt2$，则 $a^b=2$ 有理。两种情况都给出一对“无理底数/无理指数/有理幂”，无需判定 $x$ 属于哪类。
  \item 把颜色编码为 $\mathbb Z_q$。第一人公开前方颜色和的模 $q$ 值，自己的颜色可能错；第 $j$ 人减去自己看到的颜色与此前正确答案，即可反解自身颜色。协议依赖固定顺序、公共编码和无噪声广播；有噪声时应加重复/纠错码，增加的冗余取决于信道错误率和允许的失败概率。
\end{enumerate}

\subsection*{基础衔接：独立计算}
恰为 4 的倍数。从倍数只能到非倍数；从非倍数取其模 4 的余数可回到倍数，连同 $n=0$ 边界完成归纳。
